\documentclass[12pt,spanish,openany,letterpaper,pagesize]{scrbook}
\usepackage[utf8]{inputenc}
\usepackage[spanish]{babel}%escribir con acentos sin necesidad de comandos \'{} .
\usepackage{fancyhdr}
%\usepackage{epsfig}
\usepackage{epic}
\usepackage{eepic}
\usepackage{amsmath, amssymb, amsthm, latexsym, graphics, graphpap, layout, wasysym, multicol, enumerate,amsfonts,mathrsfs}
\usepackage{graphics,lscape}
\usepackage[mathscr]{euscript}
\usepackage{xcolor}
\usepackage{verbatim}
\usepackage{booktabs}
\usepackage{float}
\usepackage{hyperref}
\usepackage{threeparttable}
\usepackage{amscd}
\usepackage{here}
\usepackage[pdftex]{graphicx}
\usepackage{lscape}
\usepackage{tabularx}
\usepackage{subfigure}
\usepackage{longtable}
\usepackage{geometry}
\usepackage{cite}
\usepackage{algorithmic}
\usepackage{diagbox}
\usepackage{array,multirow}
\usepackage[all]{ xy}

\usepackage{rotating} %Para rotar texto, objetos y tablas seite. No se ve en DVI solo en PS. Seite 328 Hundebuch
                        %se usa junto con \rotate, \sidewidestable ....

\renewcommand{\theequation}{\thechapter-\arabic{equation}}
\renewcommand{\thefigure}{\textbf{\thechapter-\arabic{figure}}}
\renewcommand{\thetable}{\textbf{\thechapter-\arabic{table}}}

\pagestyle{fancyplain}%\addtolength{\headwidth}{\marginparwidth}
\textheight22.5cm \topmargin0cm \textwidth16.5cm
\oddsidemargin0.5cm \evensidemargin-0.5cm%
\renewcommand{\chaptermark}[1]{\markboth{\thechapter\; #1}{}}
\renewcommand{\sectionmark}[1]{\markright{\thesection\; #1}}
\lhead[\fancyplain{}{\thepage}]{\fancyplain{}{\rightmark}}
\rhead[\fancyplain{}{\leftmark}]{\fancyplain{}{\thepage}}
\fancyfoot{}
\thispagestyle{fancy}%

\addtolength{\headwidth}{0cm}
\unitlength1mm %Define la unidad LE para Figuras
%\mathindent0cm %Define la distancia de las formulas al texto,  fleqn las %descentra
\marginparwidth0cm
\parindent0cm %Define la distancia de la primera linea de un parrafo a la margen

%Para tablas,  redefine el backslash en tablas donde se define la posici\'{o}n del texto en las
%casillas (con \centering \raggedright o \raggedleft)
\newcommand{\PreserveBackslash}[1]{\let\temp=\\#1\let\\=\temp}
\let\PBS=\PreserveBackslash

%Espacio entre lineas
\renewcommand{\baselinestretch}{1.1}

%Neuer Befehl f\"{u}r die Tabelle Eigenschaften der Aktivkohlen
\newcommand{\arr}[1]{\raisebox{1.5ex}[0cm][0cm]{#1}}

%Neue Kommandos
\usepackage{Befehle}
%Inicio del documento. Tener en cuenta que hay archivos auxiliares

\newtheorem{teor}{Teorema}[section]
\theoremstyle{definition}
\newtheorem{definition}[teor]{Definición}
\theoremstyle{example}
\newtheorem{example}[teor]{Ejemplo}
\newtheorem{lema}[teor]{Lema}
\newtheorem{afirmacion}[teor]{Afirmación}
\newtheorem{proposicion}[teor]{Proposición}
\newtheorem{coro}[teor]{Corolario}
\newtheorem{ole}[teor]{Conjetura}
%\theoremstyle{proposicion}

%\def\proof{\paragraph{Demostración:\\}}%\begin{proof}
%\def\endproof{\hfill$\blacksquare$}%\end{proof}

\begin{document}
\pagenumbering{roman}
%\newpage
%\setcounter{page}{1}
\begin{center}
\begin{figure}
\centering%
\includegraphics[scale=0.4]{Imagenes/logo_1.png}
%{file=uc.jpg,scale=1}%
\end{figure}
\thispagestyle{empty} \vspace*{2.0cm} \textbf{\huge
Predicción del comportamiento de una enfermedad simulada en autómatas celulares con un algoritmo propuesto en redes neuronales}\\[6.0cm]
\Large\textbf{Jorge Andrés Ibáñez Huertas}\\[4.0cm]
\small Universidad Central\\
Departamento de Matemáticas\\
Bogotá, Colombia\\
2021\\
\end{center}

\newpage{\pagestyle{empty}\cleardoublepage}

\newpage
\begin{center}
\thispagestyle{empty} \vspace*{0cm} \textbf{\huge
Predicción del comportamiento de una enfermedad simulada en autómatas celulares con un algoritmo propuesto en redes neuronales}\\[3.5cm]
\Large\textbf{Jorge Andrés Ibáñez Huertas}\\[3.0cm]
\small Trabajo de grado presentado como requisito parcial para optar al
t\'{\i}tulo de:\\
\textbf{Matemático}\\[3.0cm]
Director:\\
Carlos Isaac Zainea\\[3.5cm]
Universidad Central\\
Departamento de Matemáticas\\
Bogotá, Colombia\\
2021\\
\end{center}

\newpage{\pagestyle{empty}\cleardoublepage}

% \newpage
% \thispagestyle{empty} \textbf{}\normalsize
% \\\\\\%
% \textit{"Tan sólo por la educación puede el hombre llegar a ser hombre.\\ El hombre no es más que lo que la educación hace de él"\\
% \textit{Immanuel Kant}}\\[4.0cm]

\begin{flushright}
\begin{minipage}{8cm}
    \noindent
        \small
\end{minipage}
\end{flushright}

\newpage{\pagestyle{empty}\cleardoublepage}

\newpage
\thispagestyle{empty} \textbf{}\normalsize
\\\\\\%
\textbf{\LARGE Agradecimientos}\\

% Agradezco a mis padres Roosvelt Rivera y Angelica Useche por los años de apoyo incondicional en mis años de estudio, además de la profesora Edel María Serrano Iglesias (Q.E.P.D) sin la cuál no hubiera entrado al mundo de las matemáticas.\\
% También agradezco a los profesores que a lo largo de la carrera estuvieron dispuestos a enseñar y guiar con incontable paciencia y dedicación, entre ellos: Xiomara Rojas, Diana Herrera, Diana Pulido, Isaac Zainea, Miguel Pachón, Henry Naranjo, Fabián Sánchez, Nicolás Avilán y mi director de tesis Henry Sánchez.

\newpage{\pagestyle{empty}\cleardoublepage}

\newpage
\textbf{\LARGE Resumen}\\

% En este trabajo se desarrollan los resultados obtenidos en el artículo  de Morales \& Arbieto,  específicamente la distancia $GH^0$ dada en ese manuscrito. Una vez se haya comprendido la fundamentación y  analizado algunas propiedades de dicha métrica, se procederá a estudiar la estabilidad topológica de las dinámicas sobre espacios compactos utilizando en principio la noción de Walters con una métrica $C^0$. Luego, se define y analiza la noción de $GH$-estabilidad  resaltando la relación que existe entre ambas nociones.\\

\textbf{\LARGE Abstract}\\\\

% In this work the results obtained in the article by Morales \& Arbieto are developed, specifically, the distance $ GH^0$ given in the article. Once the foundation is understood and some properties of that metric have been analyzed, we will proceed to study the topological stability of the dynamics on compact spaces using a notion given by Walters with a  $C^0$-metric. After that, will be defined and analyzed the $GH$-stability notion highlighting the relationship between both notions.

\renewcommand{\tablename}{\textbf{Tabla}}
\renewcommand{\figurename}{\textbf{Figura}}
\renewcommand{\listtablename}{Lista de Tablas}
\renewcommand{\listfigurename}{Lista de Figuras}
\renewcommand{\contentsname}{Contenido}

%\newcommand{\clearemptydoublepage}{\newpage{\pagestyle{empty}\cleardoublepage}}
\cleardoublepage
\addcontentsline{toc}{chapter}{Lista de figuras} % para que aparezca en el indice de contenidos
\listoffigures % indice de figuras

%\cleardoublepage
%\addcontentsline{toc}{chapter}{Lista de tablas} % para que aparezca en el indice de contenidos

\tableofcontents %% this produces the table of contents - you might have guessed :-)

%\listoftables % indice de tablas

%\chapter*{Lista de s\'{\i}mbolos}
\addcontentsline{toc}{chapter}{\numberline{}Lista de s\'{\i}mbolos}
Esta secci\'{o}n es opcional, dado que existen disciplinas que no manejan s\'{\i}mbolos y/o abreviaturas.\\

Se incluyen s\'{\i}mbolos generales (con letras latinas y griegas), sub\'{\i}ndices, super\'{\i}ndices y abreviaturas (incluir s\'{o}lo las clases de s\'{\i}mbolos que se utilicen). Cada una de estas listas debe estar ubicada en orden alfab\'{e}tico de acuerdo con la primera letra del s\'{\i}mbolo.
\section*{S\'{\i}mbolos con letras latinas}
 \label{simbolos}
 \renewcommand{\arraystretch}{1.3}
%\begin{longtable}[l]{*{4}{>{$}l<{$}}p{9cm}}
\begin{longtable}[l]{>{$}l<{$}l>{$}l<{$}>{$}l<{$}}
%\begin{tabular}
\textbf{S\'{\i}mbolo}&\textbf{T\'{e}rmino}&\textbf{Unidad SI}&\textbf{Definici\'{o}n}\\[0.5ex]\hline
\endfirsthead%
\textbf{S\'{\i}mbolo}&\textbf{T\'{e}rmino}&\textbf{Unidad SI}&\textbf{Definici\'{o}n}\\[0.5ex]\hline
\endhead%
      A              &\'{A}rea                                   &\text{m}^{2}                         &\int\int dxdy\\%
      A_{\text{BET}} &\'{A}rea interna del s\'{o}lido                &\frac{\text{m}^{2}}{\text{g}}        &\text{ver DIN ISO 9277}\\%
      A_{\text{g}}   &\'{A}rea transversal de la fase gaseosa    &\text{m}^{2}                         &\text{Ec...}\\%
      A_{\text{s}}   &\'{A}rea transversal de la carga a granel  &\text{m}^{2}                         &\text{Ec...}\\%
      a              &Coeficiente                            &1                                    &\text{Ec...}\\%
      a              &Contenido de ceniza                    &1                                    &\frac{m_{\text{ceniza}}}{m_{\text{bm,0}}}\\%
      c              &Contenido de carbono                   &1                                    &\frac{m_{\text{C}}}{m}\\%
      c              &Longitud de la cuerda                  &\text{m}                             &\text{Figura...}\\
      c              &Concentraci\'{o}n de la cantidad de materia&\frac{\text{mol}}{\text{m}^{3}}      &\frac{n}{V}\\%
      D              &Di\'{a}metro                               &\text{m}                             &\\%
      E_{\text{A}}   &Energ\'{\i}a de activaci\'{o}n                  &\frac{\text{kJ}}{\text{mol}}         &\text{Ec....}\\%
      F              &Fracci\'{o}n de materia vol\'{a}til            &1                                    &\text{ver DIN 51720}\\%
      Fr             &N\'{u}mero de Froude                       &1                                    &\frac{\omega^{2}R}{g_{\text{0}}}\\%
      \overrightarrow{g}&Aceleraci\'{o}n de la gravedad          &\frac{\text{m}}{\text{s}^{2}}        &\frac{d^{2}\overrightarrow{r}}{dt^{2}}\\%
      H              &Entalp\'{\i}a                               &\text{J}                             &U+PV\\%
      H_{\text{o}}   &Poder calor\'{\i}fico superior              &\frac{\text{MJ}}{\text{kg}}          &\text{ver DIN 51857}\\%
      h              &Contenido de hidr\'{o}geno                 &1                                    &\frac{m_{\text{H}}}{m}\\%
      K              &Coeficiente de equilibrio              &1                                    &\text{Ec...}\\%
      L              &Longitud                               &\text{m}                             &DF\\%
      L              &Longitud del reactor                   &\text{m}                             &\text{Figura...}\\%
      m              &Masa                                   &\text{kg}                            &DF\\%
      \dot{m}        &Flujo de masa                          &\frac{\text{kg}}{\text{s}}           &\frac{m}{t}\\%
      n              &Velocidad de rotaci\'{o}n                  &\frac{\text{1}}{\text{s}}            &\frac{\omega}{2\pi}\\%
      n              &Cantidad de materia                    &\text{mol}                           &DF\\%
      P              &Presi\'{o}n                                &\text{Pa}                            &\frac{\vec{F}\cdot\vec{n}}{A}\\%
      Q              &Calor                                  &\text{kJ}                            &\text{1. $LT$}\\%
      T              &Temperatura                            &\text{K}                             &DF\\%
      t              &Tiempo                                 &\text{s}                             &DF\\%
      x_{\text{i}}   &Fracci\'{o}n de la cantidad de materia     &1                                    &\frac{n_{\text{i}}}{n}\\%
      V              &Volumen                                &\text{m}^{3}                         &\int{dr^{3}}\\%
      \vec{u}        &Velocidad                              &\frac{\text{m}}{\text{s}}            &(\frac{dr}{dt},r\frac{d\upsilon}{dt},\frac{dz}{dt})\\%
      w_{\text{i}}   &Fracci\'{o}n en masa del componente i      &1                                    &\frac{m_{\text{i}}}{m_{\text{0}}}\\%
      w_{\text{w,i}} &Contenido de humedad de la sustancia i &1                                    &\frac{m_{\text{\wasser}}}{m_{\text{i,0}}}\\%
      Z              &Factor de gases reales                 &1                                    &\frac{pv}{RT}\\%
\end{longtable}
\vspace{5ex}
\section*{S\'{\i}mbolos con letras griegas}

\begin{longtable}[l]{>{$}l<{$}l>{$}l<{$}>{$}l<{$}}
\textbf{S\'{\i}mbolo}&\textbf{T\'{e}rmino}&\textbf{Unidad SI}&\textbf{Definici\'{o}n}\\[0.5ex] \hline%
\endfirsthead%
\textbf{S\'{\i}mbolo}&\textbf{T\'{e}rmino}&\textbf{Unidad SI}&\textbf{Definici\'{o}n}\\[0.5ex] \hline%
\endhead%
\renewcommand{\arraystretch}{1.3}
 \label{simbolosg}
     \alpha_{\text{BET}}  &Factor de superficie                  &\frac{\text{m}^{2}}{\text{g}}   &(w_{\text{F,waf}})(A_{\text{BET}})\\%
     \beta_{\text{i}}     &Grado de formaci\'{o}n del componente i   &1                               &\frac{m_{\text{i}}}{m_{\text{bm,0}}}\\%
     \gamma               &Wandhaftreibwinkel (Stahlblech)       &1                               &\text{Secci\'{o}n...}\\
     \epsilon             &Porosidad de la part\'{\i}cula             &1                               &1-\frac{\rho_{\text{s}}}{\rho_{\text{w}}}\\%
     \eta                 &mittlere Bettneigungswinkel (St\"{u}rzen) &1                               &\text{Figura...}\\%
     \theta               &\'{A}ngulo de inclinaci\'{o}n de la cama      &1                               &\text{Figura...}\\
     \theta_{\text{O}}    &\'{A}ngulo superior de avalancha          &1                               &\text{Figura...}\\
     \theta_{\text{U}}    &\'{A}ngulo inferior de avalancha          &1                               &\text{Figura...}\\
     \kappa               &Velocidad de calentamientoe           &\frac{\text{K}}{\text{s}}       &\frac{dT}{dt}\\%
     \nu                  &Coeficiente estequiom\'{e}trico           &1                               &\text{ver DIN 13345}\\%
     \rho_{\text{b}}      &Densidad a granel                     &\frac{\text{kg}}{\text{m}^{3}}  &\frac{m_{\text{S}}}{V_{\text{S}}}\;(\text{Secci\'{o}n...})\\
     \rho_{\text{s}}      &Densidad aparente                     &\frac{\text{kg}}{\text{m}^{3}}  &\frac{m_{\text{F}}}{V_{\text{P}}}\;(\text{Secci\'{o}n...})\\
     \rho_{\text{w}}      &Densidad verdadera                    &\frac{\text{kg}}{\text{m}^{3}}  &\frac{m_{\text{F}}}{V_{\text{F}}}\;(\text{Secci\'{o}n...})\\
     \tau                 &Tiempo adimensional                   &1                               &\text{Ec....}\\%
     \Phi_{\text{V}}      &Flujo volum\'{e}trico                     &\frac{\text{m}^{3}}{\text{s}}   &\frac{\Delta V}{\Delta t}\\
     \omega               &Velocidad angular                     &\frac{1}{\text{s}}              &\frac{d\upsilon}{dt}\\

\end{longtable}


\section*{Sub\'{\i}ndices}
\begin{longtable}[l]{>{}l<{}l}
  \textbf{Sub\'{\i}ndice} & \textbf{T\'{e}rmino} \\[0.5ex] \hline%
  \endfirsthead%
  \textbf{Sub\'{\i}ndice} & \textbf{T\'{e}rmino} \\[0.5ex] \hline%
  \endhead%
\renewcommand{\arraystretch}{1.4}\label{simbolosg}

 bm&materia org\'{a}nica\\%
 DR&Dubinin-Radushkevich\\%
 E&Experimental\\%
 g&Fase gaseosa\\%
 k&Condensado\\%
 Ma&Macroporos\\%
 P&Part\'{\i}cula\\%
 p&Poro\\%
 p&Pirolizado\\%
 R&Reacci\'{o}n\\%
 t&Total\\%
 wf&Libre de agua\\%
 waf&Libre de agua y de ceniza\\%
 0&Estado de referencia\\%

\end{longtable}


\setlength{\extrarowheight}{0pt}


\section*{Super\'{\i}ndices}
\begin{longtable}[l]{>{}l<{}l}
  \textbf{Super\'{\i}ndice} & \textbf{T\'{e}rmino} \\[0.5ex] \hline%
  \endfirsthead%
  \textbf{Super\'{\i}ndice} & \textbf{T\'{e}rmino} \\[0.5ex] \hline%
  \endhead%
\renewcommand{\arraystretch}{1.4}\label{simbolosg}

 n &Coeficiente x\\%



\end{longtable}


\setlength{\extrarowheight}{0pt}


\section*{Abreviaturas}
\begin{longtable}[l]{>{}l<{}l}
  \textbf{Abreviatura} & \textbf{T\'{e}rmino} \\[0.5ex] \hline%
  \endfirsthead%
  \textbf{Abreviatura} & \textbf{T\'{e}rmino} \\[0.5ex] \hline%
  \endhead%
\renewcommand{\arraystretch}{1.4}\label{simbolosg}
 1.$LT$&Primera ley de la termodin\'{a}mica\\%
 $DF$    &Dimensi\'{o}n fundamental\\%
 $RFF$   &Racimos de fruta fresca\\%

\end{longtable}


\setlength{\extrarowheight}{0pt}
%\include{Resumen}%\newcommand{\clearemptydoublepage}{\newpage{\pagestyle{empty}\cleardoublepage}}

\pagenumbering{arabic}
\include{Contenidos/C1-Epidemiología}
% \chapter{Conclusiones}

% Al realizar este trabajo se comprendió la manera en que se define la distancia $C^0$,\,la distancia de \textit{Hausdorff},\, la distancia \textit{Gromov-Hausdorff}. Así mismo, cómo se pueden calcular distancias entre objetos en donde estas métricas estén definidas usando las diversas caracterizaciones de distancia, así como propiedades que éstas o los objetos (conjuntos, espacios, funciones) poseen. Posterior a eso, fue posible entender qué motiva y cómo funciona la denominada distancia $GH^0$. Ya con estos conceptos claros, se logró comprender de qué manera difieren o son similares ciertos espacios y/o funciones en términos de qué tan cercanos son con respecto a una forma de definir esa distancia. Luego de eso, fue posible a través del entendimiento del concepto de estabilidad topológica de Walters \cite{WalPOTP}, estudiar la propuesta de $GH$-estabilidad topológica dada por Arbieto \& Morales en \cite{AM}.\\

% Las conclusiones que se pueden enunciar de este trabajo son:
% \begin{itemize}
%     % \item Mediante la introducción de métricas o pseudo métricas es posible estudiar y entender el comportamiento de espacios y sistemas, en el sentido de  comparar con elementos cercanos. A su vez, estudiar fenómenos de convergencia y estabilidad.
%     \item Todo homeomorfismo topológicamente estable es isométricamente estable. Sin embargo el recíproco es falso en general.
%     \item Un homeomorfismo topológicamente $GH$-estable es isometricamente estable. Esto quiere decir que una función $GH$-estable ante pequeñas perturbaciones no solo produce una función con una estructura dinámica, sino que además la dinámica misma de la perturbación es muy cercana. 
%     \item Existe un homeomorfismo el cual es topológicamente estable pero no es topológicamente $GH$-estable. Por ende, la $GH$-estabilidad es una condición distinta de estabilidad.
%     \item Todo homeomorfismo topológicamente $GH$-estable en el circulo es topológicamente estable.
%     \item Todo homeomorfismo  expansivo con la P.O.T.P  de un espacio métrico compacto es topológicamente $GH$-estable (extendiendo el resultado dado por Walters).
%     \item Toda función continua positivamente expansiva con la $P.O.T.P_+$ de un espacio métrico compacto es topológicamente $GH-$estable (Caso no invertible).
%     \item Lo desarrollado en este trabajo sugiere que en espacios compactos la $GH$-estabilidad es una condición más fuerte que la estabilidad topológica. Sin embargo, esto aún no ha sido probado.
% \end{itemize}
% \include{Contenidos/C3-Apéndices}

\addcontentsline{toc}{chapter}{\numberline{}Bibliografía}
\bibliographystyle{plain} % apalike
\bibliography{BibliMSc}
\end{document}